\documentclass[11pt]{letter}
\usepackage[margin=2.5cm]{geometry}
\usepackage{hyperref}
\usepackage[T1]{fontenc}

\signature{Prof.\ Dr.\ Tobias-Benedikt Blask\\
Harz University of Applied Sciences\\
Friedrichstraße 57--59\\
38855 Wernigerode, Germany\\
\texttt{tblask@hs-harz.de}}

\date{\today}

\begin{document}

\begin{letter}{%
Prof.\ Luca Mora and Prof.\ Mei-Chih Hu\\
Editors-in-Chief\\
\textit{Technological Forecasting and Social Change}\\
Elsevier}

\opening{Dear Professors Mora and Hu,}

I am pleased to submit the manuscript entitled \textbf{``When the Middle Disappears: Three Orders of AI-Driven Economic Transformation from a Multi-LLM Delphi Analysis''} for consideration as a regular research article in \textit{Technological Forecasting and Social Change}.

\textbf{Summary.} The paper applies the Multi-LLM Delphi --- a two-round foresight method in which seven frontier language models from three geopolitical regions (US, China, EU) serve as structured panelists --- to AI-driven economic transformation across eight GICS-aligned industry verticals, yielding an approximately 114,000-word corpus analyzed through structured comparative content analysis with consensus scoring. The study identifies six high-consensus convergence zones, including labor market bifurcation (7/7 models), infrastructure bottleneck repricing (6/7), and the collapse of intermediary business models (5/7). A central finding is the ``apprenticeship crash'' --- the prediction that AI's elimination of junior knowledge-work roles will sever experiential training pipelines --- which moves from 1/7 to 7/7 endorsement after controlled feedback. Geopolitical divergence persists and intensifies under feedback pressure.

\textbf{Fit with TFSC.} The manuscript addresses TFSC's core scope at the intersection of technological innovation and social change. It contributes to three streams represented in the journal: (1)~the Delphi methodology literature (von der Gracht, 2012; Ecken et al., 2011; Sokolov et al., 2025), which it extends by introducing LLMs as structured panelists; (2)~the AI-and-labor literature (Choi and Leigh, 2024; Qian et al., 2023; Grybauskas and Cardenas-Rubio, 2024), which it enriches with cross-sector, multi-order analysis; and (3)~the sociotechnical transitions literature grounded in GPT theory and the Multi-Level Perspective (Geels, 2005, 2020; Korzinov and Savin, 2018; Coccia, 2018), which it applies to AI as a general-purpose technology. 20 of the 83 references are published in TFSC.

\textbf{Novelty.} The study's primary contribution is substantive: a comprehensive cross-sector map of AI's cascading effects from task automation through structural reorganization to systemic transformation, grounded in the three-order effects framework (Hilty et al., 2006; Hall, 1993). Its secondary contribution is methodological: to the author's knowledge, this is among the first studies to systematically employ multiple frontier LLMs as structured panelists within a formal two-round Delphi protocol with consensus scoring.

\textbf{Transparency note.} The manuscript was produced using the Open Academic Paper Machine (\url{https://github.com/TobiasBlask/open-paper-machine}), an AI-assisted research workflow tool developed by the author. The author conceived, designed, and directed the entire study; AI tools were used for prose drafting and production under author orchestration. Full disclosure is provided in the AI Declaration section. I take responsibility for the content.

The manuscript has not been published previously and is not under consideration at any other journal. All data and analytical materials --- including the formal codebook, consensus scoring matrices, prompt instruments, figure generation specifications, and the complete raw model outputs from both rounds --- are available through the supplementary repository at \url{https://github.com/TobiasBlask/when-the-middle-disappears}.

\closing{Sincerely,}

\end{letter}
\end{document}
